% drafts/08-hamiltonian.tex — From Hamiltonian to Pauli Strings
%
% Target length: ~1.5 pages (two-column PRA)
% Shows the end-to-end pipeline as a concrete worked example.

The preceding sections have developed the algebraic framework in the
abstract.  We now trace the complete pipeline from a second-quantized
Hamiltonian to a sum of Pauli strings, making the computational
structure concrete.

%%====================================================================
\subsection{The encoding pipeline}
\label{sec:ham-pipeline}

Given a second-quantized Hamiltonian~\eqref{eq:hamiltonian} and an
encoding (specified by either an index-set scheme or a tree), the
transformation to Pauli strings proceeds in five stages:

\begin{enumerate}
  \item \textbf{Normal ordering.}  Each product of ladder operators
    is rewritten in normal order (all $\ad{}$ to the left of all
    $\an{}$) using the combining algebra
    (\cref{def:combining-algebra}).  This step applies the CAR
    repeatedly, potentially doubling the number of terms at each swap
    (when the $\delta_{ij}$ identity fires).  The result is a sum of
    normal-ordered products with updated coefficients.

  \item \textbf{Majorana decomposition.}  Each ladder operator is
    replaced by its Majorana decomposition~\eqref{eq:ladder-from-majorana}:
    \begin{equation*}
      \ad{j} \to \tfrac{1}{2}(c_j - id_j), \qquad
      \an{j} \to \tfrac{1}{2}(c_j + id_j).
    \end{equation*}
    Each two-operator product $\ad{p}\an{q}$ expands to four Majorana
    products; a four-operator product $\ad{p}\ad{q}\an{s}\an{r}$
    expands to sixteen.

  \item \textbf{Majorana $\to$ Pauli.}  Each Majorana operator is
    mapped to a single Pauli string via the chosen encoding's
    construction (\cref{sec:index-set} or \cref{sec:tree-encoding}).
    This is a lookup: $c_j \mapsto S_j^{(c)}$ and
    $d_j \mapsto S_j^{(d)}$.

  \item \textbf{Pauli multiplication.}  Each product of Pauli strings
    is reduced to a single Pauli string by qubit-by-qubit
    multiplication~\eqref{eq:pauli-mult}, with phase tracked
    symbolically in $\Zi$ (\cref{sec:exactness}).

  \item \textbf{Like-term combination.}  All resulting Pauli strings
    are collected in a signature-keyed dictionary.  Terms with the
    same Pauli signature have their coefficients summed; terms whose
    coefficients cancel to zero are removed.  The output is the
    final Pauli Hamiltonian~\eqref{eq:pauli-hamiltonian}.
\end{enumerate}

\paragraph{Complexity.}
For an $n$-mode molecular Hamiltonian, the one-body part has $O(n^2)$
terms and the two-body part has $O(n^4)$ terms.  Each term generates
at most 16~Majorana products (step~2), each involving at most 4~Pauli
string multiplications (step~4), each costing $O(n)$ (qubit-by-qubit).
The total cost is $O(n^5)$, matching existing
implementations~\cite{mcclean2020, bergholm2022}.  Like-term
combination (step~5) adds $O(T \log T)$ where $T$ is the total number
of intermediate Pauli strings.

%%====================================================================
\subsection{Worked example: H$_2$ in minimal basis}
\label{sec:ham-example}

The H$_2$ molecule in the STO-3G basis has four spin-orbitals
($n = 4$) and a Hamiltonian with 13 distinct ladder-operator terms
(after symmetry reduction).  We trace a single representative
two-body term through the pipeline using the Jordan--Wigner encoding.

\paragraph{Input.}
Consider the term $g = 0.1744\;\ad{0}\ad{2}\an{3}\an{1}$ from the
two-electron integral.

\paragraph{Step 1: Normal ordering.}
This term is already in normal order ($\ad{}\ad{}\an{}\an{}$), so
no swaps are needed.

\paragraph{Step 2: Majorana decomposition.}
\begin{align*}
  \ad{0}\ad{2}\an{3}\an{1}
  &= \tfrac{1}{16}\,(c_0 - id_0)(c_2 - id_2)(c_3 + id_3)(c_1 + id_1).
\end{align*}
Expanding yields 16 products of four Majorana operators, each with a
coefficient in $\frac{1}{16}\{+1, -1, +i, -i\}$.

\paragraph{Step 3: Majorana $\to$ Pauli (JW).}
Under Jordan--Wigner, the Majorana operators for $n = 4$ are:
\begin{align*}
  c_0 &= X_0,         & d_0 &= Y_0, \\
  c_1 &= Z_0 X_1,     & d_1 &= Z_0 Y_1, \\
  c_2 &= Z_0 Z_1 X_2, & d_2 &= Z_0 Z_1 Y_2, \\
  c_3 &= Z_0 Z_1 Z_2 X_3, & d_3 &= Z_0 Z_1 Z_2 Y_3.
\end{align*}

\paragraph{Step 4: Pauli multiplication.}
One of the 16 products is
$c_0 \cdot c_2 \cdot c_3 \cdot c_1
= (XIII)(ZZXI)(ZZZX)(ZXII)$.
Multiplying qubit-by-qubit with $\Zi$ phase tracking:
\begin{align*}
  \text{qubit 0:}\quad & X \cdot Z \cdot Z \cdot Z
    = X \cdot Z \cdot I = X \cdot I = X, \quad \phi_0 = (-i)(+1)(+1) = -i \\
  \text{qubit 1:}\quad & I \cdot Z \cdot Z \cdot X = I \cdot I \cdot X = X,
    \quad \phi_1 = (+1)(+1)(+1) \\
  \text{qubit 2:}\quad & I \cdot X \cdot Z \cdot I = (-iY) \cdot I = -iY,
    \quad \phi_2 = (-i)(+1) \\
  \text{qubit 3:}\quad & I \cdot I \cdot X \cdot I = X,
    \quad \phi_3 = +1.
\end{align*}
Global phase: $\phi = (-i)(-i) = -1$.
Result: $(-1) \cdot XXYXI = -XXYXI$.

After accounting for the $\frac{1}{16}$ prefactor and the original
coefficient: $0.1744 \times \frac{1}{16} \times (-1) \cdot XXYX$.

\paragraph{Step 5: Combination.}
After processing all 16 products from this term (and all 12 other
terms), the 15~surviving Pauli strings with nonzero coefficients are:
\begin{equation*}
  H_\text{qubit} = \sum_{k=1}^{15} h_k\, \sigma_k,
\end{equation*}
where the specific signatures ($IIII$, $IIIZ$, $IIZI$, $XXYY$,
$YYXX$, etc.) and coefficients match the standard H$_2$ STO-3G
benchmark.

%%====================================================================
\subsection{What the pipeline makes explicit}
\label{sec:ham-observations}

Several properties of this pipeline deserve emphasis:

\begin{enumerate}
  \item \textbf{Encoding-independence of the pipeline.}  Only step~3
    depends on the encoding choice.  Steps~1, 2, 4, and~5 are
    universal.  Changing the encoding is a one-line substitution:
    replace the Majorana lookup table.

  \item \textbf{Phase tracking is separable.}  Steps~2 and~4 each
    produce phases in $\Zi$, which accumulate via the group law
    independently of the coefficient arithmetic.  This separation is
    what makes exact phase tracking possible
    (\cref{thm:exactness}).

  \item \textbf{Like-term cancellation is essential.}  The 16
    Majorana products from a single two-body term often produce
    Pauli strings that partially cancel.  A four-mode Hamiltonian
    with 208 raw Pauli products (13 terms $\times$ 16 products) reduces
    to just 15 distinct Pauli strings after combination.  Missing
    cancellations produce a Hamiltonian with wrong eigenvalues.

  \item \textbf{The combinatorial explosion is bounded.}  Despite the
    $O(n^4)$ two-body terms each producing 16 Majorana products, the
    number of \emph{distinct} Pauli strings in the final Hamiltonian
    is at most $O(n^4)$ (bounded by the number of input terms), and
    in practice much smaller due to cancellations.  The dictionary
    structure handles this efficiently.
\end{enumerate}
